\documentclass[UTF8,fontset=fandol,11pt,a4paper]{ctexart}
\usepackage[a4paper,margin=1.75cm,headheight=15pt]{geometry}
\usepackage{amsmath,amssymb,mathtools,booktabs,tabularx,array,enumitem}
\usepackage[most]{tcolorbox}
\usepackage{xcolor,fancyhdr,hyperref,microtype,lastpage}
\newcolumntype{Y}{>{\raggedright\arraybackslash}X}
\newcolumntype{L}[1]{>{\raggedright\arraybackslash}p{#1}}
\newcolumntype{C}[1]{>{\centering\arraybackslash}p{#1}}
\setlength{\parindent}{2em}\setlength{\parskip}{0.34em}\setlength{\emergencystretch}{3em}
\renewcommand{\arraystretch}{1.2}\setlength{\tabcolsep}{3pt}
\setlist[itemize]{itemsep=0.18em,topsep=0.35em,leftmargin=2em}
\setlist[enumerate]{itemsep=0.18em,topsep=0.35em,leftmargin=2em}
\definecolor{deep}{RGB}{38,58,72}\definecolor{soft}{RGB}{244,247,248}\definecolor{greenish}{RGB}{52,91,74}\definecolor{warnc}{RGB}{136,61,58}\definecolor{purple}{RGB}{84,72,112}
\hypersetup{colorlinks=true,linkcolor=deep,urlcolor=purple,pdftitle={CO-I01 一条指令的一生}}
\pagestyle{fancy}\fancyhf{}
\lhead{\small\color{deep}\textbf{408 · CO-I01}}
\rhead{\small\color{deep}LOAD Fast/Slow Path · Precise Commit}
\cfoot{\small 第 \thepage\ 页 / 共 \pageref{LastPage} 页}\renewcommand{\headrulewidth}{0.4pt}
\newtcolorbox{corebox}[1]{breakable,colback=soft,colframe=deep,title=\textbf{#1},boxrule=0.8pt,arc=2mm}
\newtcolorbox{methodbox}[1]{breakable,colback=white,colframe=greenish,title=\textbf{#1},boxrule=0.7pt,arc=1.5mm}
\newtcolorbox{warnbox}[1]{breakable,colback=white,colframe=warnc,title=\textbf{#1},boxrule=0.7pt,arc=1.5mm}
\newcommand{\kw}[1]{\textbf{\color{deep}#1}}
\title{\vspace{-1.1cm}\textbf{CO-I01｜一条指令的一生}\\[0.4em]\large 以 LOAD 追踪语义、硬件路径、快慢分支与精确提交}
\author{408 · 计算机组成原理综合专题}
\date{}

\begin{document}
\maketitle

\begin{corebox}{完整执行流程}
\[
\boxed{\text{ISA Semantic}\to\text{Fetch / Decode}\to\text{Datapath / Pipeline}\to\text{EA}\to\text{Translation / Cache / Memory}\to\text{Result Ready}\to\text{Precise Commit}}
\]
母指令取 \texttt{LOAD rd, disp(rs)}。它既需要一次指令访问，也需要一次数据访问；任一阶段可以变慢或异常，但只有全部前置条件成立，\texttt{rd} 才能被架构提交。
\end{corebox}

\begin{methodbox}{本手册的范围}
本册只追踪一条指令跨越多个模块时的顺序、交接字段、快路径/慢路径和停止条件。位串、ISA、数据通路、流水线、存储、Cache、地址翻译和总线的具体定义仍以各自专题正文为准；本册不复制第二套定义。
\end{methodbox}

\section{术语先行：整条指令为什么需要交接}
\begin{methodbox}{五个流程词}
\begin{itemize}
\item \kw{LOAD：}从有效地址读取指定宽度数据，并把扩展后的结果写入目标寄存器的指令。
\item \kw{交接字段：}一个阶段交给下一阶段的地址、数据、权限、状态和重试信息。
\item \kw{快路径：}取指、翻译、Cache 和执行都命中且没有冲突的常规路径。
\item \kw{慢路径：}遇到缺失、等待、异常或总线争用，需要暂停、修复或重试的路径。
\item \kw{精确提交：}发生异常时，较老指令的效果已经可见，出错指令和更年轻指令的禁止副作用仍不可见。
\end{itemize}
\end{methodbox}

\tableofcontents
\newpage

\section{第一步不是画硬件，而是写 LOAD 的架构状态差}
CO-02 给出
\[
EA=Reg[rs]+SignExt(disp),\qquad value=M[EA,\,size],
\]
\[
Read=\{PC,Reg[rs],instruction\ bytes,M[EA]\},\quad
Write=\{Reg[rd],NextPC\},
\]
并附 access size、signed/unsigned extension、alignment、privilege 和可能异常。CO-01 决定 displacement 怎样扩展、内存字节怎样解释并压入目标位宽。

CO-B01 把它标准化为
\[
\{Read,Derived,Write,NextPC,Memory,Exception\},
\]
再展开为 value dependency：PC 产生取指地址；指令字产生字段；$Reg[rs]$ 与扩展立即数产生 EA；翻译和 Cache/Memory 产生 load value；只有无异常结果可进入 $Reg[rd]$。

\section{完整轨迹包含两次访问}
\subsection{Instruction fetch}
当前 PC 先按 ISA 对齐/权限规则形成 instruction access request。I-TLB/页表将 VA 翻成 PA，I-Cache 检查指令块副本；miss 时向下一级/主存取块。指令字 ready 后才可稳定译码。教学五级流水常把这一过程压缩为 IF 一拍，题设若改变 TLB/Cache latency 就必须展开。

\subsection{Decode 与 operand read}
译码恢复 opcode、寄存器字段、displacement、访问宽度和符号扩展规则；寄存器堆读取 base operand。非法指令、权限或字段约束可能在此形成异常，不能把位串成功读取等同于语义合法。

\subsection{EA generation}
ALU 计算 $EA=base+extended\ displacement$。这里产生的是 virtual effective address；alignment/check 发生在哪一层由 ISA/实现给出。EA ready 不等于 load data ready，也不等于可以写回。

\section{数据访问：翻译与 Cache 的双重身份证明}
CO-B02 接收 $(VA,access\ type,privilege,size)$，翻译输出 $(PA,permission,attributes,status)$。合法 PA 再进入 Cache 的 tag/index/offset 身份检查：
\[
\text{translation valid}\land\text{permission OK}\land\text{cache identity hit}
\Rightarrow\text{data may return}.
\]
PIPT 通常先翻译后索引；VIPT 可用 page-offset 内的 index 与 TLB 并行，但最终仍用 PA tag 和权限确认。TLB miss、Page Fault 与 Cache miss 必须分别路由。

\section{Fast Path：所有局部预测都成功}
在题设五级 in-order、I/D TLB hit、I/D Cache hit、无结构冲突和异常时，可写成：
\begin{center}\small
\begin{tabularx}{\linewidth}{L{1.4cm}L{2.3cm}YY}\toprule
阶段 & 主要结果 & 负责的硬件层 & 交给下一步\\\midrule
IF & instruction bytes ready & CO04 + CO07/06 & instruction packet\\
ID & semantic/dependency packet & CO02 + CO-B01 & operands/control intent\\
EX & EA ready & CO01/03 & data access request\\
MEM & PA valid, Cache hit, data ready & CO07 + CO-B02 + CO06 & typed load value\\
WB/Commit & $Reg[rd]$、NextPC 可见 & CO03/04 & next architectural state\\
\bottomrule
\end{tabularx}
\end{center}
阶段名不是机制证明；题目若在 EX 前移/后移分支判断、让 memory 多周期或改变寄存器堆时序，必须用 need/ready 重新排拍。

\section{Slow Path：谁暂停，谁修复，哪里重试}
\begin{center}\small
\begin{tabularx}{\linewidth}{L{2.1cm}L{2.2cm}YY}\toprule
事件 & 缺失/冲突 & 处理者 & 对指令的影响\\\midrule
I-TLB/I-Cache miss & 指令翻译/块副本 & 翻译或存储层次 & IF 等待/重试\\
RAW hazard & producer value 尚未到 need point & forwarding/interlock & forward 或 stall\\
D-TLB miss & 数据翻译副本 & page walker/refill & 数据访问重试\\
Page/permission fault & 映射/权限失败 & 精确异常 + OS & 年轻指令作废，修复后 restart 或终止\\
D-Cache miss & 数据块副本缺失 & Cache/memory controller & victim/fill/retry，load value 延后\\
Memory/bus busy & 共享资源未获权 & arbiter/controller & transaction 等待，不等于 ISR\\
\bottomrule
\end{tabularx}
\end{center}

Page Fault 进入操作系统后超出本册的硬件流程；本册只保留 fault 的交接信息与重启边界。若访问 MMIO，地址属性可能绕过普通 Cache，由 I/O 控制器和总线事务接管；普通 DRAM fill 则按主存组织和访问成本完成。

\section{Ready、Forward 与 Commit 必须分开}
load value 从 Cache/memory 返回是 \kw{ready}；旁路可把它交给年轻 consumer，是 \kw{forwardable}；只有该 LOAD 成为允许改变软件可见状态的指令时才 \kw{commit}。经典 load-use hazard 来自 consumer need 早于 load ready，而不是寄存器名字相同本身。

精确异常要求：发生 fault 的 LOAD 不写 $Reg[rd]$，更年轻指令不留下架构副作用，更老指令已按模型完成；修复后从 ISA 规定的 restart point 重新执行。内部可曾做过推测访问，不等于软件可见提交。

\begin{warnbox}{三条常见错误链}
EA 已算出 $\Rightarrow$ 数据已取回；TLB miss $\Rightarrow$ Page Fault；Cache 返回数据 $\Rightarrow$ LOAD 已提交。三条都跳过了中间 handoff 与状态条件。
\end{warnbox}

\section{性能只在成本边界正确时组合}
总 CPU time 仍是 $IC\times CPI\times T_{clk}$。单条 LOAD 的额外等待可能来自 hazard、TLB、Cache、memory 和 bus，但不能把已包含在 miss penalty 中的下一级 AMAT 再加一次，也不能同时把可并行的 TLB/Cache latency 串行求和。

先写事件树和条件概率，再根据题设重叠关系组合。例如近似数据访问成本可写
\[
T_D=T_{TLB/cache\ fast}+P(TLB\ miss)T_{walk,extra}
+P(Cache\ miss\mid translation\ succeeds)T_{fill,extra},
\]
其中 page fault 通常应作为独立稀有 slow path，且其 OS/磁盘成本不能混入普通 Cache miss。

\section{母轨迹怎样迁移到其他指令}
\begin{center}\small
\begin{tabularx}{\linewidth}{L{1.6cm}YYY}\toprule
指令 & 保留 & 删除/替换 & 提交条件\\\midrule
ADD & fetch/decode/execute/timing & 无数据翻译与 Cache；ALU 产值 & 无异常且结果到 commit\\
STORE & fetch/EA/translation/Cache & 无 $Reg[rd]$ 写回；写缓冲/策略接管 & 内存副作用按架构顺序可见\\
BRANCH & fetch/decode/compare/NextPC & 无 data load；加入预测/flush & 正确路径 PC 与年轻指令状态一致\\
MMIO LOAD & fetch/EA/translation & 普通 Cache 路径可能被属性旁路 & 控制器事务合法完成且无 fault\\
\bottomrule
\end{tabularx}
\end{center}
迁移方法是删改 handoff，不是重画一套全新 CPU。若某指令要求新数据源、目的或异常边界，先回 CO-02/B01 修改语义包，再判断 CO-03 通路是否支持。

\section{五列概念边界}
\begin{center}\small
\begin{tabularx}{\linewidth}{L{1.8cm}C{0.35cm}L{1.9cm}YY}\toprule
概念 A & $\neq$ & 概念 B & 真正区别与题面信号 & 混淆后果\\\midrule
EA ready & $\neq$ & load value ready & 地址已生成 vs 数据已返回 & 漏翻译/Cache\\
TLB miss & $\neq$ & Page Fault & 翻译副本缺失 vs 映射/权限失败 & 错写 OS 介入\\
Cache hit & $\neq$ & permission OK & 数据副本存在 vs 访问合法 & 非法数据提交\\
forward & $\neq$ & commit & 内部传值 vs 架构状态可见 & 精确异常破坏\\
bus wait & $\neq$ & interrupt & 共享事务等待 vs CPU 控制流事件 & 把停顿写成保存现场\\
综合流程 & $\neq$ & 单一专题 & 跨模块过程 vs 局部机制定义 & 产生第二套定义\\
\bottomrule
\end{tabularx}
\end{center}

\section{综合题调用协议}
\begin{enumerate}
\item 写 ISA state delta 与异常集合；
\item 把 instruction fetch 与 data access 分开；
\item 逐值标 producer、ready、consumer need 和 commit；
\item 对 VA 访问写 translation result，再做 Cache identity；
\item 每个 miss/fault/冲突写检测者、处理者、重试点；
\item 标共享资源、允许重叠项和成本是否已包含；
\item 在无异常路径证明 $Reg[rd]$ 与 NextPC 的精确提交；
\item 局部推导不足时回到对应专题的定义，不在本册补第二套定义。
\end{enumerate}

\begin{methodbox}{压缩成第一动作}
看到“分析一条指令全过程”，先画三条线：\kw{值依赖线、时间 ready/need 线、异常/提交线}；再把每个节点归入对应专题。
\end{methodbox}

\section{各阶段的机制入口与停止条件}
\begin{center}\small
\begin{tabularx}{\linewidth}{L{2.5cm}YY}\toprule
问题信号 & 进入的专题 & 本综合流程的停止点\\\midrule
\shortstack[l]{位宽 / 扩展\\舍入} & CO-01 & 得到 typed finite-width value\\
\shortstack[l]{编码 / EA\\异常语义} & CO-02 & 得到 state delta packet\\
\shortstack[l]{MUX / 微操作\\控制字} & CO-03 & 得到单指令路径与控制\\
\shortstack[l]{hazard / stall\\flush / commit} & CO-04 & 得到 need-ready 与提交顺序\\
\shortstack[l]{DRAM / 模块\\介质成本} & CO-05 & 得到访问服务模型\\
\shortstack[l]{tag / miss\\write policy} & CO-06 & 得到副本状态转换\\
\shortstack[l]{TLB / PTE\\page walk} & CO-07 & 得到 translation result\\
\shortstack[l]{bus / MMIO\\DMA} & CO-08 & 得到 transaction completion\\
两侧 handoff & CO-B01/B02 & 得到标准接口包\\
\bottomrule
\end{tabularx}
\end{center}

\begin{corebox}{一句话}
一条 LOAD 的一生不是固定五阶段口诀，而是“语义承诺沿值依赖和时间路径逐层获得证据，任何 slow path 都必须回到明确重试点，最终只在无异常且顺序合法时提交”。
\end{corebox}

\end{document}
